\scanpage{9}
\mspage{7}
\uline{Cor. 2} Soit $I\in S$ tel que $\delta I\subset L_r$.
Alors $\exists J\in S$ tel que $\tilde I\subset\tilde J^\circ$.
\margin{\deleted{\unsure{extension}}\inserted{\unsure{prolongement}} d'un segment \deleted{\unsure{avec}} à extrémités régulières}

$F_5)$ L'ordre d'un lieu régulier (on sait qu'il est $\geq2$) est égal à deux.

Équivalent à

$F_5')$ Soit $x\in L_r$, et considérons les $\tilde I^\circ$ ($I\in S$) tels que $x\in\tilde I^\circ$.
Ceux-ci forment une base de filtre, en d'autres termes, si $I',I''\in S$ tels que $x\in\tilde I'^\circ\cap\tilde I''^\circ$, $\exists I\in S$ tels que
\[
x\in\tilde I^\circ\subset\tilde I'^\circ\cap\tilde I''^\circ
\]
(NB \unsure{DPS} $=\quad\tilde I\subset\tilde I'^\circ\cap\tilde I''^\circ$).

\margin{Axiome des lieux réguliers ou de l'ordre 2}
\margin{(7.6.) Si on considère l'ens des points (= ``lieux'') d'un espace \unsure{$\sigma$} modéré, en prenant comme ``segments'' les parties \uline{modérées} qui sont isom. à $[0,1]$, alors conditions $F_1):F_4)$ satisfaites, mais \unsure{$F_5)$ (ou $F_6$ plus bas\ldots)}}

\deleted{
\uline{Remarque} Soit $S_0\deleted{\unsure{C}}\subset S$, et soit
\[
L'=\bigcup_{I\in S_0}\tilde I,\qquad S'=\{I\in S\mid\tilde I\subset L'\}
\]
(donc $S'\supset S_0$). On a deux
\[
\begin{array}{lll}
S'\unsure{-|\longrightarrow}\mathcal P(L')&&I\unsure{\mapstochar\relbar|\mapsto}\tilde I\\
S'\unsure{-|\longrightarrow}\mathcal P_2(L')&&I\unsure{\mapstochar\relbar|\mapsto}\delta I
\end{array}
\]
Les axiomes $F_1):F_5)$ sont-ils satisfaits?
}
\mspage{8}

\rotatebox{90}{\unsure{EXTRA STRONG}}
\rotatebox{90}{\unsure{EXTRA STRONG}}
